\cotefile{134-2}
\scanpage{95}
\mspage{\unsure{\ldots)|2}}
had been done and the homotopy category expressed in terms of $\infty$-groupoids. But finally it appears this is not so, they have been working throughout with a notion of $\infty$-groupoid too restrictive for the purposes I had in mind (probably because they insist I guess on strict associativity of compositions, rather than associativity up to a (given) isomorphism\inserted{\handwritten{or rather, homotopy}}) - to the effect that the simply connected homotopy types they obtain are merely products of Eilenberg-MacLane spaces, too bad ! They do not seem to have realized yet that this makes their set-up wholly unadequate \handwritten{\unsure{as}} a sweeping foundational set up for homotopy. This brings to the fore again to work out the suitable definitions for $n$-groupoids - if this is not done yet anywhere. I spent the afternoon today trying to figure out a reasonable definition, to get a feeling at least of where the difficulties are, if any. I am guided mainly of course by the topological interpretation. It will be short enough to say how far I got. The main part of the structure it seems is expressed by the sets $F_i$ ($i\in\underline{N}$) of $i$-objects, the source, target and identity maps
\[
\begin{aligned}
s^i_1,t^i_1 &: F_i\longrightarrow F_{i-1} &&(i\geq 1)\\
k^i_1 &: F_i\longrightarrow F_{i+1} &&(i\in\underline{N})
\end{aligned}
\]
\margin{\handwritten{\uline{NB} we write also $\mathrm{id}_a$ instead of ${\unsure{K|k}}^i_1(a)$ and $\check{u}$ for $\mathrm{inv}_i(u)$}}
and the symmetry map (passage to the inverse)
\[
\mathrm{inv}_i:\quad F_i\longrightarrow F_i\qquad(i\geq 1)
\]
satisfying some obvious relations : $k^i_1$ is right inverse to the \deleted{target} source and target maps $s_1^{i\deleted{\unsure{-}}+1}$, $t_1^{i+1}$, $\mathrm{inv}_i$ is an involution and ``exchanges'' source and target, and moreover for $i\geq 2$
\[
\begin{aligned}
s^{i-1}_1s^i_1&=s^{\deleted{\unsure{i-1}}i-1}_1t^i_1 &&(\mathrel{\overset{\mathrm{def}}{=}}s^i_2:F_i\longrightarrow F_{i-2})\\
t^{i-1}_1s^i_1&=t^{i-1}_1t^i_1 &&(\mathrel{\overset{\mathrm{def}}{=}}t^i_2:F_i\longrightarrow F_{i-2})\quad;
\end{aligned}
\]
\unsure{thus|thur} the composition of the source and target maps yields, for $0\leq j\leq i$, just \uline{two} maps
\[
s^i_\ell,\ t^i_\ell:\ F_i\longrightarrow F_{i-\ell}\handwritten{=F_j}\quad(\ell=i-j)\,.
\]
The next basic structure is the composition structure, where the usual composition of arrows, more specifically of $i$-objects $(i\geq1)$ $v\circ u$ (defined when $t_1(u)=s_1(v)$) must be supplemented by the Godement-type operations $\handwritten{\mu}*\handwritten{\lambda}$ when $\handwritten{\mu}$ and $\handwritten{\lambda}$ are ``arrows between arrows'' \inserted{etc} . Following this line of thought, one \unsure{getc|gets} the composition maps
\[
(u,v)\longmapsto v\overset{\unsure{i|t}}{*}_\ell u\quad:\quad (F_i,\deleted{s}t^i_\ell)\mathbin{\mathrm{x}}_{F_{i-\ell}}(F_i,s^i_\ell)\longrightarrow F_i\ ,
\]
the $\deleted{\unsure{*_\ell}}$ \deleted{\unsure{application|operation}}\inserted{composition of} $i$-objects \inserted{$u,v$} $(1\leq\ell\leq i)$ being defined when the $\ell$-target of $u$ is equal to the $\ell$-source of $v$, and then we have
\[
\left.\begin{aligned}
\deleted{\unsure{*_\ell}}\ s^i_1(v*_\ell u)&=s^i_1(v)\underset{\unsure{1}}{*_{\ell-1}}s^i_1(u)\\
\handwritten{t^i_1(v*_\ell u)}&\handwritten{{}=t^i_1(v)*_{\ell-1}t^i_1(u)}
\end{aligned}\right\}\handwritten{\quad\ell\geq2\quad\text{i.e.}\quad\ell-1\geq1}
\]
% archive sheet number: 94
